\chapter{Ephemeris and Almanac Design}
\label{chap:ephemeris-almanac}

This chapter specifies the mathematical contract of \lupnt{}'s broadcast
\emph{navigation-message} models for lunar satellites: the precise,
short-validity \cpp{LansEphemeris} and the coarse, long-validity
\cpp{LansAlmanac}. Both live in
\file{cpp/lupnt/applications/ephemeris/}: the models in
\file{cpp/lupnt/applications/ephemeris/lunanet_ephemeris.cc} and
\file{cpp/lupnt/applications/ephemeris/lunanet_almanac.cc}, their shared
numerics in \file{cpp/lupnt/applications/ephemeris/ephemeris_basis.h}, the
bit-budget analysis in
\file{cpp/lupnt/applications/ephemeris/ephemeris_app.cc}, and the runtime
message generator in
\file{cpp/lupnt/applications/ephemeris/ephemeris_gen_app.cc}.

The chapter fixes, in order: the state, frame and unit contract shared by both
models (\cref{sec:eph-contract}); the ephemeris representation --- an
osculating two-body Kepler baseline plus a Chebyshev residual plus an optional
argument-of-latitude Fourier residual (\cref{sec:eph-representation}); the
least-squares fit that determines it (\cref{sec:eph-fit}); the almanac
element-wise series representation, its mean-anomaly residual and
argument-of-latitude correction (\cref{sec:eph-almanac,sec:eph-almanac-fit});
the bit allocation and quantization that size a message against the LunaNet
budget (\cref{sec:eph-bits}); the runtime broadcast generator
(\cref{sec:eph-broadcast}); and the model boundaries
(\cref{sec:eph-boundaries}).

The algorithms follow \citep[Ch.~7, Ch.~8]{iiyama2026thesis} --- Chapter~7 for
the ephemeris, Chapter~8 for the almanac --- and the companion papers
\citep{iiyama2025ephemeris,cortinovis2023ephemeris,cortinovis2024navi}. The
bit-allocation conventions mirror the GPS broadcast interface
\citep{icdgps200}, and the message budget is the one allotted by the LunaNet
Interoperability Specification \citep{lunanet2022}. Where the code
deliberately departs from the thesis formulation, the difference is recorded
in a note rather than silently harmonised.

% ===========================================================================
\section{State, Frame, and Unit Contract}
\label{sec:eph-contract}

\subsection{Sampled trajectory and the three entry points}

Both models represent a sampled Cartesian trajectory
\begin{equation}
  \label{eq:eph-state-contract}
  \mat{RV} \in \mathbb{R}^{N\times 6},
  \qquad
  \mat{RV}_{i,:} = \big[\vec{r}(t_i)\transpose \;\; \vec{v}(t_i)\transpose\big],
  \qquad
  \vec{r}\ \text{in \si{\meter}},\quad \vec{v}\ \text{in \si{\meter\per\second}},
\end{equation}
sampled at epochs $t_i$ given in seconds, strictly increasing, and measured
from an arbitrary but fixed origin. Nothing in either model depends on the
absolute epoch: only the differences $t_k = t - t_\mathrm{ref}$ enter the
basis functions, so the caller may use elapsed simulation seconds, seconds
past J2000 in $\TDB$, or any other consistent scale. The gravitational
parameter $GM$ supplied through the options must be consistent with those
units (\SI{4.902800118e12}{\meter\cubed\per\second\squared} for the Moon, the
default \code{GM\_MOON}).

Three entry points are shared by \cpp{LansEphemeris} and \cpp{LansAlmanac}:

\begin{itemize}[leftmargin=1.4em]
  \item \cpp{Fit(t\_s, rv)} returns a flat parameter vector $\vec{\alpha}$
    whose layout is documented element-by-element by \cpp{ParamNames()} and
    whose length is \cpp{NumParams()}.
  \item \cpp{Eval(t\_s, params)} reconstructs $\mat{RV}$ at arbitrary query
    epochs from $\vec{\alpha}$ alone. This is the receiver-side operation: it
    must not need anything that is not broadcast.
  \item \cpp{EvalError(t\_s, rv\_ref, params)} returns an
    \cpp{EphemerisFitErrorStats}, the RTN-decomposed RMS and 95th-percentile
    position and velocity errors against a reference arc.
\end{itemize}

The error statistics are computed by \cpp{ComputeFitErrorStats} in
\file{cpp/lupnt/applications/ephemeris/ephemeris_basis.h}. The difference
$\delta\mat{RV} = \mat{RV}_\mathrm{fit} - \mat{RV}_\mathrm{ref}$ is rotated
into the Radial--Transverse--Normal triad built at each \emph{reference}
sample,
\begin{equation}
  \label{eq:eph-rtn}
  \mat{M}_\mathrm{RTN}(\vec{r},\vec{v}) =
  \begin{bmatrix} \uvec{R}\transpose \\ \uvec{T}\transpose \\ \uvec{N}\transpose \end{bmatrix},
  \qquad
  \uvec{R} = \frac{\vec{r}}{\lVert\vec{r}\rVert},\quad
  \uvec{N} = \frac{\vec{r}\times\vec{v}}{\lVert\vec{r}\times\vec{v}\rVert},\quad
  \uvec{T} = \uvec{N}\times\uvec{R},
\end{equation}
and the four reported components are $[R, T, N, \text{3-D norm}]$. The
95th percentile uses the nearest-rank definition, index
$\lceil 0.95N\rceil - 1$ of the sorted absolute values, so it is exactly
reproducible and never interpolates.

\subsection{The three frames}

Three frames appear in the specification \citep[\S7.1.1]{iiyama2026thesis};
they are summarised in \cref{tab:eph-frames}.

\begin{table}[H]
  \centering
  \caption{The three frames used by the navigation-message models. Only MCI and
    PA are \lupnt{} \cpp{Frame} enumerators; PAI is an internal construction of
    the fitting code and is never named in the public API.}
  \label{tab:eph-frames}
  \begin{tabularx}{\textwidth}{@{}l l L@{}}
    \toprule
    Frame & Enumerator & Role \\
    \midrule
    MCI & \code{Frame::MOON\_CI} &
      Moon-centred inertial, axes aligned with the ICRF. The default fitting
      frame; the two-body baseline is a standard osculating Kepler orbit and
      the frame spin rate is zero. \\
    PA & \code{Frame::MOON\_PA} &
      Moon-fixed principal-axis frame, rotating about its $z$-axis at
      $\omega_b = \Omega_\mathrm{Moon} = \SI{2.6617e-6}{\radian\per\second}$
      (\code{OMEGA\_MOON}). The frame in which the broadcast message is
      normally expressed. \\
    PAI & --- (internal) &
      Principal-Axis Inertial: the instantaneous orientation of PA frozen at
      each epoch. A quasi-inertial frame used only as the intermediate in
      which osculating elements are computed and interpreted. \\
    \bottomrule
  \end{tabularx}
\end{table}

Neither model ever calls \cpp{ConvertFrame}. Instead the entire PA/PAI
relationship is carried by a single scalar --- the frame spin rate --- and a
single vector operation, the $\vec{\omega}\times\vec{r}$ cross product with
$\vec{\omega} = \omega_b\uvec{z}$:
\begin{equation}
  \label{eq:eph-spin}
  \omega_b(\mathcal{F}) =
  \begin{cases}
    \Omega_\mathrm{Moon}, & \mathcal{F} = \code{MOON\_PA},\\[2pt]
    0, & \text{otherwise},
  \end{cases}
  \qquad
  \vec{\omega}\times\vec{r} = \omega_b
  \begin{bmatrix} -r_y \\ r_x \\ 0 \end{bmatrix}.
\end{equation}

\begin{lstlisting}[language=cpp, caption={Frame spin rate and the $z$-axis cross product, \code{cpp/lupnt/applications/ephemeris/lunanet\_ephemeris.cc}. The same two helpers are repeated verbatim in \code{lunanet\_almanac.cc}.}, label={lst:eph-spin}]
// cpp/lupnt/applications/ephemeris/lunanet_ephemeris.cc :: FrameSpinRate / SpinCross

// Angular velocity [rad/s] of `frame` about its z-axis relative to inertial:
// OMEGA_MOON for the rotating Moon-fixed principal-axis frame, 0 otherwise
// (so an inertial `frame`, e.g. MOON_CI, reproduces a standard Kepler orbit).
double FrameSpinRate(Frame frame) { return frame == Frame::MOON_PA ? OMEGA_MOON : 0.0; }

// z-axis (Moon pole) cross product omega x r, with omega = spin * z_hat.
Vec3d SpinCross(double spin, const Vec3d& r) { return spin * Vec3d(-r(1), r(0), 0.0); }
\end{lstlisting}

Setting $\omega_b = 0$ collapses every PA-specific term, so the inertial case
is not a separate code path but the $\omega_b \to 0$ limit of the rotating
one. When \code{options.frame == Frame::MOON\_PA} the code therefore
\emph{fits in PAI and returns in PA}: the caller supplies MOON\_PA states, the
fit adds the $\vec{\omega}\times\vec{r}$ offset to recover the PAI velocity,
\begin{equation}
  \label{eq:eph-pai-velocity}
  \vec{r}^{\,\mathrm{PAI}} = \vec{r}^{\,\mathrm{PA}},
  \qquad
  \vec{v}^{\,\mathrm{PAI}} = \vec{v}^{\,\mathrm{PA}} + \vec{\omega}\times\vec{r}^{\,\mathrm{PA}},
\end{equation}
and \cpp{Eval} subtracts it back on output. Position is untouched: PA and PAI
share their axes instantaneously, and differ only in the transport term that
appears in the velocity.

\begin{lupntnote}
  The frame handling is what lets a modest polynomial stay accurate. In
  MOON\_PA the ascending node of the fitted two-body baseline drifts at
  $-\omega_b$, so the baseline's node tracks the Moon's rotation instead of
  being chased by the polynomial residual. \cpp{EphemerisGenApp} therefore
  broadcasts with \code{output\_frame = MOON\_PA}
  (\file{cpp/lupnt/applications/ephemeris/ephemeris_gen_app.cc}), and the
  shipped study configuration \file{configs/ephemeris.yaml} sets
  \code{output\_frame: MOON\_PA} while propagating truth in MOON\_CI.
\end{lupntnote}

% ===========================================================================
\section{Ephemeris Representation}
\label{sec:eph-representation}

\cpp{LansEphemeris} \citep[\S7.2, Algorithm~3]{iiyama2026thesis} models
Cartesian position as an osculating two-body Kepler baseline plus a Chebyshev
polynomial residual and an optional Fourier residual in the argument of
latitude, with velocity obtained from the \emph{analytic} time derivatives of
all three parts. Nothing is finite-differenced.

\subsection{Base parameters}

The eight base parameters are the reference epoch, the fit span, and the six
osculating classical elements evaluated once, in PAI, at the window midpoint
$t_0$:
\begin{equation}
  \label{eq:eph-base-params}
  \vec{\alpha}_\mathrm{base} =
  \big[\,t_\mathrm{ref},\ t_\mathrm{fit},\ a,\ e,\ i,\ \Omega,\ \omega,\ M_\mathrm{ref}\,\big]\transpose .
\end{equation}
Here $t_\mathrm{ref} = t_{\lfloor N/2\rfloor}$ is the midpoint sample (not the
arc mean), $t_\mathrm{fit} = t_{N-1} - t_0$ is the full span, and
$(a,e,i,\Omega,\omega,M_\mathrm{ref})$ are \cpp{CartToClassical} applied to the
midpoint state after the PAI velocity correction \cref{eq:eph-pai-velocity}.
Local time is measured from the midpoint,
\begin{equation}
  \label{eq:eph-tk}
  t_k = t - t_\mathrm{ref} \in \Big[-\tfrac{t_\mathrm{fit}}{2},\ \tfrac{t_\mathrm{fit}}{2}\Big],
  \qquad
  z = \frac{2t_k}{t_\mathrm{fit}} \in [-1, 1],
\end{equation}
so that $z$ is the normalized Chebyshev argument. Choosing the midpoint rather
than the start of the arc halves the magnitude of $t_k$ and centres the
polynomial, which is why the elements are taken there.

If \code{use\_keplerian\_baseline} is \code{false} the six element slots are
dropped and only $[t_\mathrm{ref}, t_\mathrm{fit}]$ remain; the model then
degenerates to a pure Cartesian polynomial, useful only as a baseline for
comparison because it needs a far higher order for the same accuracy.

\subsection{Chebyshev basis and its analytic derivative}

The Chebyshev polynomials of the first kind, constant term first, are
generated by the standard three-term recurrence
\citep[Eqs.~7.1--7.2]{iiyama2026thesis}
\begin{equation}
  \label{eq:eph-cheb-rec}
  T_0(z) = 1,
  \qquad
  T_1(z) = z,
  \qquad
  T_j(z) = 2z\,T_{j-1}(z) - T_{j-2}(z),
  \quad j = 2,\dots,n_c .
\end{equation}
Differentiating \cref{eq:eph-cheb-rec} with respect to time, and using
$\dd z/\dd t = 2/t_\mathrm{fit}$ (a constant, because $t_\mathrm{fit}$ is
fixed for the window), gives a recurrence for the derivative basis that runs
alongside the value recurrence:
\begin{equation}
  \label{eq:eph-cheb-dt}
  \dot T_0 = 0,
  \qquad
  \dot T_1 = \frac{2}{t_\mathrm{fit}},
  \qquad
  \dot T_j = 2z\,\dot T_{j-1} + \frac{4}{t_\mathrm{fit}}\,T_{j-1} - \dot T_{j-2}.
\end{equation}
This is the product rule applied term by term: the $2z T_{j-1}$ product
contributes $2\dot z T_{j-1} + 2z\dot T_{j-1}$, and $2\dot z = 4/t_\mathrm{fit}$.
Because $\dot T_j$ is exact rather than a divided difference, the broadcast
velocity carries no step-size error and no additional parameters --- the same
coefficients $\vec{c}$ serve both position and velocity.

Both recurrences are evaluated column-wise over all $N$ samples at once,
returning $N\times(n_c+1)$ matrices.

\begin{lstlisting}[language=cpp, caption={Chebyshev value and derivative bases, \code{cpp/lupnt/applications/ephemeris/ephemeris\_basis.h}.}, label={lst:eph-cheb}]
// cpp/lupnt/applications/ephemeris/ephemeris_basis.h :: ChebyshevBasis / ChebyshevBasisDt

inline MatXd ChebyshevBasis(const VecXd& t_k, double t_fit, int order) {
  const int n = static_cast<int>(t_k.size());
  MatXd T(n, order + 1);
  T.col(0).setOnes();
  if (order == 0) return T;
  VecXd z = (2.0 / t_fit) * t_k;
  T.col(1) = z;
  for (int j = 2; j <= order; ++j) {
    T.col(j) = 2.0 * z.array() * T.col(j - 1).array() - T.col(j - 2).array();
  }
  return T;
}

inline MatXd ChebyshevBasisDt(const VecXd& t_k, double t_fit, int order) {
  const int n = static_cast<int>(t_k.size());
  MatXd T = ChebyshevBasis(t_k, t_fit, order);
  MatXd Tdot(n, order + 1);
  Tdot.col(0).setZero();
  if (order == 0) return Tdot;
  const double dzdt = 2.0 / t_fit;
  VecXd z = (2.0 / t_fit) * t_k;
  Tdot.col(1).setConstant(dzdt);
  for (int j = 2; j <= order; ++j) {
    Tdot.col(j) = 2.0 * z.array() * Tdot.col(j - 1).array() + 2.0 * T.col(j - 1).array() * dzdt
                  - Tdot.col(j - 2).array();
  }
  return Tdot;
}
\end{lstlisting}

\subsection{The osculating two-body baseline}

\cpp{EvalKeplerianBaseline} propagates the frozen elements
\cref{eq:eph-base-params} with the two-body mean motion
\begin{equation}
  \label{eq:eph-mean-motion}
  n = \sqrt{\frac{GM}{a^3}},
  \qquad
  M(t_k) = M_\mathrm{ref} + n\,t_k,
\end{equation}
converts $(a,e,i,\Omega,\omega,M(t_k))$ to a Cartesian PAI state with
\cpp{ClassicalToCart}, and --- when the output frame rotates --- rigidly
rotates that state into it. With $\theta(t_k) = -\omega_b t_k$ and
$\mat{C}(\theta) = \mat{R}_z(\theta)$,
\begin{equation}
  \label{eq:eph-baseline-rot}
  \vec{r}^{\,\mathrm{PA}}_\mathrm{OE} = \mat{C}(\theta)\,\vec{r}^{\,\mathrm{PAI}},
  \qquad
  \vec{v}^{\,\mathrm{PA}}_\mathrm{OE} = \mat{C}(\theta)\,\vec{v}^{\,\mathrm{PAI}}
    + \dot{\mat{C}}(\theta)\,\vec{r}^{\,\mathrm{PAI}},
\end{equation}
where
\begin{equation}
  \label{eq:eph-rotmat}
  \mat{C} =
  \begin{bmatrix}
    \cos\theta & -\sin\theta & 0 \\
    \sin\theta & \cos\theta & 0 \\
    0 & 0 & 1
  \end{bmatrix},
  \qquad
  \dot{\mat{C}} = \frac{\partial\mat{C}}{\partial\theta}\,\dot\theta
  = \omega_b
  \begin{bmatrix}
    \sin\theta & \cos\theta & 0 \\
    -\cos\theta & \sin\theta & 0 \\
    0 & 0 & 0
  \end{bmatrix}.
\end{equation}
Equations \cref{eq:eph-baseline-rot,eq:eph-rotmat} are Algorithm~3 steps~5--7
of \citep[Ch.~7]{iiyama2026thesis}. The net effect on the elements is that the
right ascension of the ascending node of the baseline drifts at $-\omega_b$,
which is precisely the secular signature a Moon-fixed observer sees, and the
velocity acquires the $-\vec{\omega}\times\vec{r}$ transport term. Both are
represented exactly by the baseline, so the polynomial residual never has to
absorb them.

For the Fourier terms the baseline additionally returns the argument of
latitude and its rate. With $E$ the eccentric anomaly from Newton iteration on
Kepler's equation $E - e\sin E = M$,
\begin{equation}
  \label{eq:eph-arglat}
  \nu = 2\atantwo\!\Big(\sqrt{1+e}\,\sin\tfrac{E}{2},\ \sqrt{1-e}\,\cos\tfrac{E}{2}\Big),
  \qquad
  u = \omega + \nu,
  \qquad
  \dot u = \dot\nu = \frac{n\sqrt{1-e^2}}{(1 - e\cos E)^2},
\end{equation}
the last equality following from $\dot E = n/(1-e\cos E)$ and
$\dot\nu = \dot E\sqrt{1-e^2}/(1-e\cos E)$. The half-angle $\atantwo$ form of
$\nu$ is used rather than the $\cos\nu$ form because it is single-valued over
the full revolution and stays accurate at $e = 0.6$, the eccentricity of the
ELFO reference orbit.

\subsection{Residual model and reconstruction}

The residual added to the baseline is, per Cartesian axis
$\xi \in \{x, y, z\}$,
\begin{equation}
  \label{eq:eph-residual}
  \Delta\xi^\mathrm{res}(t)
  = \sum_{j=0}^{n_c} C^{\xi}_{T,j}\, T_j\!\Big(\frac{2t_k}{t_\mathrm{fit}}\Big)
  + \sum_{h=1}^{M}\Big[C^{\xi}_{c,h}\cos(h u) + C^{\xi}_{s,h}\sin(h u)\Big],
\end{equation}
with the corresponding analytic derivative
\begin{equation}
  \label{eq:eph-residual-dot}
  \Delta\dot\xi^\mathrm{res}(t)
  = \sum_{j=0}^{n_c} C^{\xi}_{T,j}\,\dot T_j
  + \sum_{h=1}^{M} h\,\dot u\Big[-C^{\xi}_{c,h}\sin(h u) + C^{\xi}_{s,h}\cos(h u)\Big].
\end{equation}
\cpp{Eval} then reconstructs
\begin{equation}
  \label{eq:eph-reconstruct}
  \xi^\mathrm{PA} = \xi^\mathrm{PA}_\mathrm{OE} + \Delta\xi^\mathrm{res},
  \qquad
  \dot\xi^\mathrm{PA} = \dot\xi^\mathrm{PA}_\mathrm{OE} + \Delta\dot\xi^\mathrm{res},
\end{equation}
which is Algorithm~3 steps~8--9. The Fourier harmonics are functions of $u$,
not of $t$: they are phase-locked to the orbit rather than to the clock, which
is what lets a single harmonic pair capture a once-per-revolution residual
even when the revolution period is not commensurate with the fit window.

\subsection{Parameter layout}

The full flat vector produced by \cpp{Fit}, named element-by-element by
\cpp{ParamNames()}, is laid out as in \cref{tab:eph-params}. Its length is
\begin{equation}
  \label{eq:eph-numparams}
  N_\mathrm{eph} = b + 3(n_c + 1) + 6M,
  \qquad
  b = \begin{cases} 8, & \text{Keplerian baseline},\\ 2, & \text{otherwise},\end{cases}
\end{equation}
so the default $n_c = 8$, $M = 0$ configuration carries
$8 + 27 = 35$ parameters and the shipped study configuration
($n_c = 6$, $M = 0$) carries $8 + 21 = 29$.

\begin{table}[H]
  \centering
  \caption{\cpp{LansEphemeris} parameter layout. Index ranges assume the
    Keplerian baseline ($b = 8$); the Chebyshev block is grouped by axis, and
    the Fourier block follows it in full.}
  \label{tab:eph-params}
  \begin{tabularx}{\textwidth}{@{}l l l L@{}}
    \toprule
    Index & \cpp{ParamNames()} & Symbol & Meaning \\
    \midrule
    $0$ & \code{t\_ref} & $t_\mathrm{ref}$ & Midpoint sample epoch \si{\second} \\
    $1$ & \code{t\_fit} & $t_\mathrm{fit}$ & Fit span $t_{N-1}-t_0$ \si{\second} \\
    $2$ & \code{a} & $a$ & Osculating semi-major axis at $t_\mathrm{ref}$ (PAI) \\
    $3$ & \code{e} & $e$ & Eccentricity \\
    $4$ & \code{i} & $i$ & Inclination \\
    $5$ & \code{raan} & $\Omega$ & Right ascension of the ascending node \\
    $6$ & \code{argp} & $\omega$ & Argument of periapsis \\
    $7$ & \code{M\_ref} & $M_\mathrm{ref}$ & Mean anomaly at $t_\mathrm{ref}$ \\
    \midrule
    $8+d(n_c{+}1)+j$ & \code{x\_j}, \code{y\_j}, \code{z\_j} & $C^{\xi}_{T,j}$ &
      Chebyshev coefficient $j = 0,\dots,n_c$ of axis $d = 0,1,2$ \\
    $8+3(n_c{+}1)+2dM+2(h{-}1)$ & \code{x\_fc\_h}, \dots & $C^{\xi}_{c,h}$ &
      Cosine coefficient of harmonic $h = 1,\dots,M$ \\
    $8+3(n_c{+}1)+2dM+2(h{-}1)+1$ & \code{x\_fs\_h}, \dots & $C^{\xi}_{s,h}$ &
      Sine coefficient of harmonic $h$ \\
    \bottomrule
  \end{tabularx}
\end{table}

\Cref{tab:eph-opts} lists the configuration knobs.

\begin{table}[H]
  \centering
  \caption{\cpp{EphemerisFitOptions}, \code{cpp/lupnt/applications/ephemeris/lunanet\_ephemeris.h}.}
  \label{tab:eph-opts}
  \begin{tabularx}{\textwidth}{@{}l l L@{}}
    \toprule
    Field & Default & Effect \\
    \midrule
    \code{poly\_order} & $8$ &
      Chebyshev order $n_c$ of the position residual; the derivative basis of
      the same order supplies the velocity residual. \\
    \code{use\_keplerian\_baseline} & \code{true} &
      Subtract the osculating two-body orbit before fitting. Disabling it
      leaves a pure polynomial model needing much higher order. \\
    \code{gm} & \code{GM\_MOON} &
      Central-body $GM$ in \si{\meter\cubed\per\second\squared}, consistent
      with the state units. \\
    \code{frame} & \code{Frame::MOON\_CI} &
      Frame of the input states \emph{and} of the evaluated output; sets
      $\omega_b$ through \cref{eq:eph-spin}. \\
    \code{num\_fourier\_terms} & $0$ &
      Number of argument-of-latitude harmonics $M$; $0$ gives a pure
      Chebyshev $+$ Kepler model. Requires the Keplerian baseline. \\
    \bottomrule
  \end{tabularx}
\end{table}

\begin{lupntnote}
  \textbf{Implementation vs.\ thesis.} Thesis Eq.~7.3 shows a single
  second-harmonic Fourier term $C_c\cos 2u + C_s\sin 2u$. The code instead
  uses fundamental-and-multiples $\cos(hu),\sin(hu)$ for $h = 1,\dots,M$, of
  which the thesis form is the special case $M = 2$ with the first harmonic
  constrained to zero. With the default \code{num\_fourier\_terms = 0} the
  model is pure Chebyshev $+$ Kepler. Fourier terms require
  \code{use\_keplerian\_baseline}, because the argument of latitude is
  produced by the baseline orbit; \cpp{Fit} enforces this with a
  \cpp{LUPNT\_CHECK}.
\end{lupntnote}

% ===========================================================================
\section{Ephemeris Least-Squares Fitting}
\label{sec:eph-fit}

\cpp{LansEphemeris::Fit} \citep[\S7.3.1, Eqs.~7.5--7.7]{iiyama2026thesis}
proceeds in three stages: fix the osculating elements from the midpoint state,
subtract the baseline arc to form a position residual, and solve one linear
least-squares problem per axis over a shared design matrix.

\subsection{Stage 1 --- elements from the midpoint state}

The midpoint index is $\lfloor N/2\rfloor$. The state there is lifted to PAI
by \cref{eq:eph-pai-velocity} and converted with \cpp{CartToClassical}. No
iteration and no optimisation is involved: the elements are \emph{fixed}, not
estimated.

\subsection{Stage 2 --- the residual}

With the baseline arc $\vec{r}_\mathrm{OE}(t_m)$ evaluated at every sample by
\cref{eq:eph-baseline-rot}, the residual is
\begin{equation}
  \label{eq:eph-residual-def}
  \Delta\vec{r}(t_m) = \vec{r}(t_m) - \vec{r}_\mathrm{OE}(t_m),
  \qquad m = 1,\dots,N .
\end{equation}
Only \emph{position} enters the objective; the velocity is not fit at all. It
follows from \cref{eq:eph-reconstruct} as the analytic derivative of whatever
position model the fit produces, which guarantees that the broadcast position
and velocity are exactly consistent --- a receiver differentiating the
broadcast position numerically would recover the broadcast velocity.

\subsection{Stage 3 --- the design matrix and the QR solve}

Stacking the Chebyshev and Fourier bases side by side gives the
$N\times(n_c+1+2M)$ design matrix
\begin{equation}
  \label{eq:eph-design}
  \mat{T} =
  \big[\,\underbrace{\mat{T}_\mathrm{cheb}}_{N\times(n_c+1)}\ \big|\
        \underbrace{\mat{F}}_{N\times 2M}\,\big]
  =
  \begin{bmatrix}
    T_0(z_1) & \cdots & T_{n_c}(z_1) & \cos u_1 & \sin u_1 & \cdots & \cos(Mu_1) & \sin(Mu_1)\\
    T_0(z_2) & \cdots & T_{n_c}(z_2) & \cos u_2 & \sin u_2 & \cdots & \cos(Mu_2) & \sin(Mu_2)\\
    \vdots   &        & \vdots       & \vdots   & \vdots   &        & \vdots     & \vdots \\
    T_0(z_N) & \cdots & T_{n_c}(z_N) & \cos u_N & \sin u_N & \cdots & \cos(Mu_N) & \sin(Mu_N)
  \end{bmatrix},
\end{equation}
with $z_m = 2t_{k,m}/t_\mathrm{fit}$ and $u_m = u(t_{k,m})$ from
\cref{eq:eph-arglat}. The three axes are solved independently against the same
$\mat{T}$:
\begin{equation}
  \label{eq:eph-ls}
  \vec{\alpha}_d^{\,\star}
  = \argmin_{\vec{\alpha}\in\mathbb{R}^{n_c+1+2M}}
    \sum_{m=1}^{N}\big\lVert \Delta r_d(t_m) - \mat{T}_{m,:}\,\vec{\alpha} \big\rVert_2^2,
  \qquad d \in \{x, y, z\},
\end{equation}
whose solution splits as
$\vec{\alpha}_d^{\,\star} = [\,\vec{C}^{d}_{T}\;;\;\vec{C}^{d}_{F}\,]$, the
head being the $n_c+1$ Chebyshev coefficients and the tail the $2M$ Fourier
coefficients of that axis.

Because $\mat{T}$ does not depend on the axis, its column-pivoted Householder
QR factorization $\mat{T}\mat{P} = \mat{Q}\mat{R}$ is computed \emph{once} and
reused for all three right-hand sides. Column pivoting is not cosmetic: with
$M > 0$ the Chebyshev block and the Fourier block are not orthogonal to one
another, and with a short window relative to the orbital period the low
harmonics are nearly representable by the low-order polynomials. The revealed
rank and the pivoting keep the solve well behaved in that regime, where a
normal-equations solve would lose roughly half the available digits.

The sample-count guard is $N \ge n_c + 2M + 2$, so the system is never
underdetermined.

\begin{lstlisting}[language=cpp, caption={Design-matrix assembly and the shared per-axis QR solve, \code{cpp/lupnt/applications/ephemeris/lunanet\_ephemeris.cc}.}, label={lst:eph-fit}]
// cpp/lupnt/applications/ephemeris/lunanet_ephemeris.cc :: LansEphemeris::Fit

const int cheb_len = order + 1;
const int four_len = 2 * num_fourier;
const MatXd T = ChebyshevBasis(t_k, t_fit, order);
const MatXd F = num_fourier > 0 ? FourierBasis(arc.u, num_fourier) : MatXd(n, 0);

// Joint least-squares solve over [Chebyshev | Fourier] for each axis.
MatXd design(n, cheb_len + four_len);
design.leftCols(cheb_len) = T;
if (four_len > 0) design.rightCols(four_len) = F;
const auto qr = design.colPivHouseholderQr();
MatXd cheb_coeffs(cheb_len, 3);
MatXd four_coeffs(four_len, 3);
for (int d = 0; d < 3; ++d) {
  const VecXd sol = qr.solve(residual_pos.col(d));
  cheb_coeffs.col(d) = sol.head(cheb_len);
  if (four_len > 0) four_coeffs.col(d) = sol.tail(four_len);
}
\end{lstlisting}

\begin{lupntnote}
  \textbf{Implementation vs.\ thesis.} Thesis Eq.~7.4 samples the arc at
  Chebyshev--Lobatto nodes to suppress Runge oscillation near the window
  edges; the \lupnt{} model fits whatever epochs the caller supplies.
  \cpp{EphemerisGenApp::GenerateFromAgent} samples the predicted arc
  \emph{uniformly} across the validity window
  (\code{ephemeris\_fit\_samples = 121} by default), so node placement is a
  caller responsibility and is not enforced by \cpp{Fit}. Separately, the
  osculating elements and the polynomial coefficients are fit
  \emph{sequentially} --- elements first from the midpoint state, then a
  single linear solve for the residual --- matching the thesis's fast and
  robust scheme rather than a joint nonlinear solve over all
  $N_\mathrm{eph}$ parameters.
\end{lupntnote}

% ===========================================================================
\section{Almanac Representation}
\label{sec:eph-almanac}

\cpp{LansAlmanac} \citep[\S8.1, Eq.~8.1]{iiyama2026thesis}, Algorithm~2 of
\citep{iiyama2025ephemeris}, takes the opposite approach from the ephemeris:
instead of correcting a Cartesian baseline, it fits \emph{each osculating
element} directly as a low-order polynomial plus an element-specific Fourier
term over a multi-day window. The representation is therefore always
element-based, and there is no separate ``baseline'' to enable or disable.

\subsection{Element series}

For an element $\xi(t)$ in normalized time $s = t_k/t_\mathrm{fit}$ with
$t_k = t - t_\mathrm{ref}$ and $t_\mathrm{ref} = t_0$ (the \emph{first}
sample, unlike the ephemeris's midpoint),
\begin{equation}
  \label{eq:eph-alm-series}
  \xi(t) = \sum_{p=0}^{P}\beta^{\xi}_p\,s^p
    + \sum_{h=1}^{M}\Big[C^{\xi}_{c,h}\cos(h\,\omega_\xi t_k)
                        + C^{\xi}_{s,h}\sin(h\,\omega_\xi t_k)\Big].
\end{equation}
Note the deliberate asymmetry: the polynomial is in the dimensionless $s$, so
its coefficients have the units of $\xi$ regardless of window length, while
the Fourier arguments are in raw seconds $t_k$, so that the harmonic
frequencies are physical and independent of the window length.

The base frequency is element-specific --- the orbital mean motion for the
semi-major axis, the doubled sidereal harmonic for everything else
\citep[Eq.~8.1]{iiyama2026thesis}:
\begin{equation}
  \label{eq:eph-alm-freq}
  \omega_a = \sqrt{\frac{GM}{a_\mathrm{ref}^3}} = \frac{2\pi}{T_\mathrm{orb}},
  \qquad
  \omega_{e} = \omega_{i} = \omega_{\Omega} = \omega_{M} = \omega_{u}
    = \frac{4\pi}{T_\mathrm{sid}},
  \qquad T_\mathrm{sid} = \SI{27.321661}{\dayunit}.
\end{equation}
The semi-major axis oscillates predominantly at the orbital period, driven by
the short-period gravitational terms; the remaining elements are dominated by
the semi-monthly Earth-perturbation signature, hence the $4\pi/T_\mathrm{sid}$
(twice per sidereal month) base frequency. $a_\mathrm{ref}$ is the arithmetic
mean of the osculating semi-major axis over the fit window and is broadcast so
that the receiver can reproduce $\omega_a$ exactly.

\begin{lstlisting}[language=cpp, caption={Element-specific base frequency and the element design matrix, \code{cpp/lupnt/applications/ephemeris/lunanet\_almanac.cc}.}, label={lst:eph-alm-basis}]
// cpp/lupnt/applications/ephemeris/lunanet_almanac.cc :: ElementFreq / ElementBasis

// Base Fourier angular frequency [rad/s] for element `d`: the orbital mean
// motion (2*pi/T_orb) for the semi-major axis (d==0), else 4*pi/T_sid.
double ElementFreq(int d, double a_ref, double gm, double sidereal_period_s) {
  if (d == 0) return std::sqrt(gm / (a_ref * a_ref * a_ref));
  return 4.0 * PI / sidereal_period_s;
}

// Design matrix [N x SegmentLen]: 1, s, ..., s^p (s = t_k/t_fit) followed by
// cos(h*freq*t_k), sin(h*freq*t_k) for h = 1..num_fourier.
MatXd ElementBasis(const VecXd& t_k, double t_fit, int poly_order, int num_fourier,
                   double freq) {
  const int n = static_cast<int>(t_k.size());
  MatXd A(n, SegmentLen(poly_order, num_fourier));
  const VecXd s = t_k.array() / t_fit;
  A.col(0).setOnes();
  for (int p = 1; p <= poly_order; ++p) A.col(p) = A.col(p - 1).array() * s.array();
  const int off = poly_order + 1;
  for (int h = 1; h <= num_fourier; ++h) {
    A.col(off + 2 * (h - 1)) = (h * freq * t_k.array()).cos();
    A.col(off + 2 * (h - 1) + 1) = (h * freq * t_k.array()).sin();
  }
  return A;
}
\end{lstlisting}

Because only two distinct frequencies occur, only two design matrices are ever
built: $\mat{A}_a$ at $\omega_a$ for the semi-major axis, and $\mat{A}_s$ at
$4\pi/T_\mathrm{sid}$ shared by the other five elements.

\subsection{The six represented elements}

Six element series are carried: $a$, $e$, $i$, $\Omega$, the \emph{mean-anomaly
residual} $\Delta M$, and the \emph{argument-of-latitude correction}
$u_\mathrm{corr}$. The last two are the two substantive departures from a
classical almanac.

\paragraph{Mean-anomaly residual.} A polynomial cannot economically represent
the mean anomaly itself, which sweeps many multiples of $2\pi$ across a
15-day window. The almanac therefore broadcasts only the \emph{residual} about
the nominal two-body drift implied by the fitted semi-major axis
\citep[Algorithm~4, steps~4--5]{iiyama2026thesis}:
\begin{equation}
  \label{eq:eph-alm-mean}
  M(t) = \underbrace{\int_{0}^{t_k} n(\tau)\,\dd\tau}_{M_\mathrm{nom}(t)} + \Delta M(t),
  \qquad
  n(\tau) = \sqrt{\frac{GM}{a(\tau)^3}},
\end{equation}
where $a(\tau)$ is the \emph{fitted} semi-major axis series
\cref{eq:eph-alm-series}, not a constant. The integral is evaluated by
cumulative trapezoid quadrature (\cpp{CumTrapz},
\file{cpp/lupnt/applications/ephemeris/ephemeris_basis.h}) on the query grid,
\begin{equation}
  \label{eq:eph-alm-cumtrapz}
  M_\mathrm{nom}(t_1) = 0,
  \qquad
  M_\mathrm{nom}(t_i) = M_\mathrm{nom}(t_{i-1})
    + \tfrac{1}{2}\big[n(t_i) + n(t_{i-1})\big]\,(t_i - t_{i-1}),
\end{equation}
so it is exactly the same operation at fit time and at evaluation time. This
introduces a query-grid dependence, discussed in \cref{sec:eph-boundaries}.

\paragraph{Argument-of-latitude correction.} The argument of periapsis
$\omega$ is \emph{replaced} by $u = \nu + u_\mathrm{corr}$. For a near-circular
orbit $\omega$ is ill-defined; for the eccentric ELFO reference orbit
($e = 0.6$) $\omega$ is well-defined but the reconstruction of position from
$(\omega, M)$ is badly conditioned near periapsis, where a small mean-anomaly
error maps into a large along-track position error. Parameterizing by the
argument of latitude moves the fit onto the quantity the geometry actually
depends on, and keeps the reconstruction well-conditioned at linear polynomial
order.

\begin{lstlisting}[language=cpp, caption={Mean-anomaly reconstruction inside the almanac series evaluation, \code{cpp/lupnt/applications/ephemeris/lunanet\_almanac.cc}.}, label={lst:eph-alm-eval}]
// cpp/lupnt/applications/ephemeris/lunanet_almanac.cc :: EvalSeries

AlmanacSeries s;
s.a   = A_a * params.segment(kNumBaseParams + 0 * seg, seg);
s.e   = A_s * params.segment(kNumBaseParams + 1 * seg, seg);
s.inc = A_s * params.segment(kNumBaseParams + 2 * seg, seg);
s.node = A_s * params.segment(kNumBaseParams + 3 * seg, seg);
const VecXd m_res = A_s * params.segment(kNumBaseParams + 4 * seg, seg);
s.u_corr = A_s * params.segment(kNumBaseParams + 5 * seg, seg);
const VecXd n_series = (opt.gm * s.a.array().pow(-3)).sqrt();
s.m = CumTrapz(n_series, t_k) + m_res;   // nominal drift + residual
\end{lstlisting}

\subsection{Parameter layout and the 24 fitted parameters}

Each element occupies a contiguous segment of length
\begin{equation}
  \label{eq:eph-alm-seg}
  L_\mathrm{seg} = (P + 1) + 2M,
\end{equation}
and the full vector is three bookkeeping entries followed by six such
segments,
\begin{equation}
  \label{eq:eph-alm-params}
  \vec{\alpha}_\mathrm{alm} =
  \big[\,t_\mathrm{ref},\ t_\mathrm{fit},\ a_\mathrm{ref}\ \big|\
        \vec{\alpha}_a\ \big|\ \vec{\alpha}_e\ \big|\ \vec{\alpha}_i\ \big|\
        \vec{\alpha}_\Omega\ \big|\ \vec{\alpha}_{\Delta M}\ \big|\
        \vec{\alpha}_{u}\,\big]\transpose,
  \qquad
  N_\mathrm{alm} = 3 + 6L_\mathrm{seg}.
\end{equation}

With the defaults \code{poly\_order = 1} and \code{num\_fourier\_terms = 1},
$L_\mathrm{seg} = 4$ and each element carries
$[\beta_0, \beta_1, C_c, C_s]$ --- the \textbf{24 fitted parameters} of
\citep[\S8.1]{iiyama2026thesis}, plus the three bookkeeping entries, for
$N_\mathrm{alm} = 27$ slots in total. For comparison, a GPS almanac page
carries 13 parameters \citep{icdgps200}; the extra parameters buy
representation of the strongly perturbed lunar dynamics over a 15-day window,
which no 13-parameter Keplerian set can follow.

\begin{table}[H]
  \centering
  \caption{\cpp{LansAlmanac} parameter layout with the defaults $P = 1$,
    $M = 1$ ($L_\mathrm{seg} = 4$). The names are those returned by
    \cpp{ParamNames()}.}
  \label{tab:eph-alm-params}
  \begin{tabularx}{\textwidth}{@{}l l l L@{}}
    \toprule
    Index & Name & Symbol & Meaning \\
    \midrule
    $0$ & \code{t\_ref} & $t_\mathrm{ref}$ & First sample epoch \si{\second} \\
    $1$ & \code{t\_fit} & $t_\mathrm{fit}$ & Fit span \si{\second} \\
    $2$ & \code{a\_ref} & $a_\mathrm{ref}$ & Window-mean semi-major axis, sets $\omega_a$ \\
    \midrule
    $3$--$6$   & \code{a\_p0}\dots\code{a\_fs1}    & $\vec{\alpha}_a$ & Semi-major axis, frequency $\omega_a$ \\
    $7$--$10$  & \code{e\_p0}\dots\code{e\_fs1}    & $\vec{\alpha}_e$ & Eccentricity \\
    $11$--$14$ & \code{i\_p0}\dots\code{i\_fs1}    & $\vec{\alpha}_i$ & Inclination \\
    $15$--$18$ & \code{raan\_p0}\dots\code{raan\_fs1} & $\vec{\alpha}_\Omega$ & Node, fit directly (drifts secularly in PA) \\
    $19$--$22$ & \code{M\_p0}\dots\code{M\_fs1}    & $\vec{\alpha}_{\Delta M}$ & Mean-anomaly residual about \cref{eq:eph-alm-mean} \\
    $23$--$26$ & \code{u\_p0}\dots\code{u\_fs1}    & $\vec{\alpha}_u$ & Argument-of-latitude correction $u_\mathrm{corr}$ \\
    \bottomrule
  \end{tabularx}
\end{table}

\begin{table}[H]
  \centering
  \caption{\cpp{AlmanacFitOptions}, \code{cpp/lupnt/applications/ephemeris/lunanet\_almanac.h}.}
  \label{tab:eph-alm-opts}
  \begin{tabularx}{\textwidth}{@{}l l L@{}}
    \toprule
    Field & Default & Effect \\
    \midrule
    \code{poly\_order} & $1$ &
      Secular polynomial degree $P$ in $s = t_k/t_\mathrm{fit}$. Linear is
      preferred: the $u$ parameterization stays well-conditioned at $P = 1$,
      and every extra degree costs six broadcast fields. \\
    \code{num\_fourier\_terms} & $1$ &
      Harmonics $M$ per element, at the element-specific base frequency
      \cref{eq:eph-alm-freq}. \\
    \code{gm} & \code{GM\_MOON} & Central-body $GM$. \\
    \code{sidereal\_period\_s} & $27.321661\times\code{SECS\_DAY}$ &
      $T_\mathrm{sid}$, setting $4\pi/T_\mathrm{sid}$ for the five
      non-semi-major-axis elements. \\
    \code{frame} & \code{Frame::MOON\_CI} &
      Frame of the input states and of the output; sets $\omega_b$. \\
    \bottomrule
  \end{tabularx}
\end{table}

% ===========================================================================
\section{Almanac Least-Squares Fitting}
\label{sec:eph-almanac-fit}

\cpp{LansAlmanac::Fit} \citep[\S8.2, Eqs.~8.2--8.6]{iiyama2026thesis} computes
unwrapped osculating elements from the (PAI) states, builds the two design
matrices, and solves each element by QR least squares.

\subsection{Osculating elements and unwrapping}

Every sample is lifted to PAI by \cref{eq:eph-pai-velocity} and converted with
\cpp{CartToClassical}, giving an $N\times 6$ table
$[a, e, i, \Omega, \omega, M]$. The four angular columns $i, \Omega, \omega,
M$ are then \emph{unwrapped}: consecutive samples are forced to differ by less
than $\pi$ by adding multiples of $2\pi$ (\cpp{UnwrapAngles},
\file{cpp/lupnt/applications/ephemeris/ephemeris_basis.h}). Without this step
the least-squares target would contain $2\pi$ discontinuities that a smooth
basis cannot represent and that would contaminate every coefficient. The
reference semi-major axis follows as
\begin{equation}
  \label{eq:eph-alm-aref}
  a_\mathrm{ref} = \frac{1}{N}\sum_{m=1}^{N} a(t_m).
\end{equation}

\subsection{Per-element solves}

Each element is solved independently against its design matrix,
\begin{equation}
  \label{eq:eph-alm-ls}
  \vec{\alpha}_\xi^{\,\star}
  = \argmin_{\vec{\alpha}_\xi \in \mathbb{R}^{L_\mathrm{seg}}}
    \sum_{m=1}^{N}\big\lVert \xi(t_m) - \mat{A}_{\xi,m,:}\,\vec{\alpha}_\xi \big\rVert_2^2,
  \qquad
  \mat{A}_\xi \in \{\mat{A}_a, \mat{A}_s\},
\end{equation}
again with a column-pivoted Householder QR. As in \cref{sec:eph-fit} the two
factorizations $\mat{A}_a = \mat{Q}_a\mat{R}_a\mat{P}_a\transpose$ and
$\mat{A}_s = \mat{Q}_s\mat{R}_s\mat{P}_s\transpose$ are formed once and reused
across all the right-hand sides that share them. Here the pivoting matters for
a different reason: over a window shorter than $T_\mathrm{sid}/2$ the
harmonic columns of $\mat{A}_s$ become nearly collinear with the linear
polynomial column, and column pivoting keeps the split between ``secular'' and
``periodic'' from being resolved by round-off.

The node $\Omega$ is fit directly rather than through a rate-plus-offset
parameterization: in the PA frame it drifts secularly at $-\omega_b$, which the
linear polynomial term absorbs exactly.

\subsection{The mean-anomaly residual target}

The nominal drift is built from the \emph{already fitted} semi-major axis
series, so it is exactly what the receiver will reconstruct:
\begin{equation}
  \label{eq:eph-alm-mnom}
  a_\mathrm{fit} = \mat{A}_a\vec{\alpha}_a^{\,\star},
  \qquad
  n_\mathrm{fit} = \sqrt{GM\,a_\mathrm{fit}^{-3}},
  \qquad
  M_\mathrm{nom} = \operatorname{CumTrapz}(n_\mathrm{fit}, t_k),
\end{equation}
and the residual target is the unwrapped osculating mean anomaly minus that
drift, $\Delta M_\mathrm{target} = M_\mathrm{osc} - M_\mathrm{nom}$. This
sequential structure --- $a$ first, then $M$ against the $a$-implied drift ---
is what keeps $\Delta M$ small enough for a linear polynomial plus one
harmonic to represent it.

\subsection{The argument-of-latitude target}

The $u$ correction is fit against a target constructed so that it contains no
$2\pi$ jumps. With $\nu_\mathrm{osc}$ the true anomaly implied by the
\emph{osculating} $(M, e)$ pair and $\nu_\mathrm{fit}$ the true anomaly implied
by the \emph{fitted} $(M_\mathrm{fit}, e_\mathrm{fit})$ pair,
\begin{equation}
  \label{eq:eph-alm-utarget}
  u_\mathrm{target}(t_m)
  = \omega_\mathrm{osc}(t_m)
    + \atantwo\!\big(\sin(\nu_\mathrm{osc} - \nu_\mathrm{fit}),\
                     \cos(\nu_\mathrm{osc} - \nu_\mathrm{fit})\big),
\end{equation}
where $\omega_\mathrm{osc}$ is the already-unwrapped argument of periapsis and
the $\atantwo$ wraps the small true-anomaly discrepancy into $(-\pi, \pi]$.
The sum is therefore a smooth, slowly varying function: the unwrapped
$\omega$ carries the secular behaviour, and the wrapped $\Delta\nu$ carries a
small correction. Fitting $\omega + \Delta\nu$ directly, or wrapping the whole
sum, would reintroduce the discontinuities the parameterization exists to
avoid.

\begin{lstlisting}[language=cpp, caption={Mean-anomaly residual and argument-of-latitude targets, \code{cpp/lupnt/applications/ephemeris/lunanet\_almanac.cc}.}, label={lst:eph-alm-fit}]
// cpp/lupnt/applications/ephemeris/lunanet_almanac.cc :: LansAlmanac::Fit

const VecXd a_series = A_a * a_coeffs;
const VecXd e_series = A_s * e_coeffs;
const VecXd n_series = (options_.gm * a_series.array().pow(-3)).sqrt();
const VecXd m_nom = CumTrapz(n_series, t_k);
const VecXd m_coeffs = qr_s.solve(coe.col(5) - m_nom);
const VecXd m_fit = m_nom + A_s * m_coeffs;

// Argument-of-latitude correction u - nu(from fitted M), formed smoothly as
// argp + wrap(nu_osc - nu_fit) so the least-squares target has no 2*pi jumps.
VecXd u_target(n);
for (int i = 0; i < n; ++i) {
  const double nu_o = TrueAnomaly(SolveKepler(coe(i, 5), coe(i, 1)), coe(i, 1));
  const double nu_f = TrueAnomaly(SolveKepler(m_fit(i), e_series(i)), e_series(i));
  const double dnu = std::atan2(std::sin(nu_o - nu_f), std::cos(nu_o - nu_f));
  u_target(i) = coe(i, 4) + dnu;  // argp (unwrapped) + small residual
}
const VecXd u_coeffs = qr_s.solve(u_target);
\end{lstlisting}

\subsection{Reconstruction}

\cpp{Eval} evaluates the six series at the query epochs, wraps the (large,
unwrapped) mean anomaly into $(-\pi, \pi]$ so that the Kepler solver inside
\cpp{ClassicalToCart} stays well-conditioned, and calls \cpp{ClassicalToCart}
with $\code{argp} = u_\mathrm{corr}$:
\begin{equation}
  \label{eq:eph-alm-reconstruct}
  \vec{rv}^{\,\mathrm{PAI}} = \operatorname{ClassicalToCart}
    \big(a, e, i, \Omega, u_\mathrm{corr}, \operatorname{wrap}(M)\big),
  \qquad
  \vec{v}^{\,\mathrm{PA}} = \vec{v}^{\,\mathrm{PAI}} - \vec{\omega}\times\vec{r}.
\end{equation}
Because $u = \nu + u_\mathrm{corr}$ by construction, passing $u_\mathrm{corr}$
in the argument-of-periapsis slot reproduces exactly the $u$-based
reconstruction of \citep[Algorithm~4, steps~6--10]{iiyama2026thesis} while
also yielding the analytic two-body velocity for free.

\begin{lupntnote}
  \textbf{Implementation vs.\ thesis.} Thesis Algorithm~4 step~11 computes
  velocity by finite differencing the reconstructed position. The code instead
  returns the \emph{analytic two-body velocity} from \cpp{ClassicalToCart}
  (minus the frame $\vec{\omega}\times\vec{r}$ term), which avoids choosing a
  finite-difference step size and removes the associated truncation and
  cancellation errors. The class documentation comment in
  \file{cpp/lupnt/applications/ephemeris/lunanet_almanac.h} still describes the
  finite-difference formulation; the implementation in
  \file{lunanet_almanac.cc} is authoritative.
\end{lupntnote}

% ===========================================================================
\section{Message Size and Bit Budget}
\label{sec:eph-bits}

The broadcast cost is sized \emph{per parameter}, so that quantization keeps
the reconstructed position error under a tolerance
\citep[\S7.3.2, Eqs.~7.8--7.12]{iiyama2026thesis}. The implementation is in
\file{cpp/lupnt/applications/ephemeris/ephemeris_app.cc} --- the
\cpp{EphemerisApp} bit-budget helpers \cpp{SearchParamBits},
\cpp{RangeBits}, and \cpp{SummarizeWindows} --- and is deliberately kept out
of the model classes, which store and evaluate unquantized doubles.

\subsection{Fractional-bit search}

For each parameter index $j$, \cpp{SearchParamBits} binary-searches the
minimum number of fractional bits $k_j$ such that a $2^{-k}$ step in that
parameter alone changes the \emph{worst-case} position over the window by less
than the tolerance $\delta_x$:
\begin{equation}
  \label{eq:eph-bits-k}
  k_j = \min\Big\{\,k \in \mathbb{N} \;:\;
    \max_{m}\big\lVert f(\vec{\alpha} + 2^{-k}\uvec{e}_j,\ t_m)
                     - f(\vec{\alpha},\ t_m)\big\rVert_2 \le \delta_x \Big\},
\end{equation}
where $f(\cdot,t)$ is the model's position output at $t$ and $\uvec{e}_j$ is
the $j$-th unit vector. Both signs of the perturbation are evaluated and the
larger displacement is taken, because the map from parameter to position is
not symmetric --- for an eccentric orbit a $+\epsilon$ and a $-\epsilon$ step
in, say, the mean anomaly produce different peak excursions.

Using the maximum over the window rather than the RMS is a deliberate choice.
For the ELFO reference orbit a coarse step in a period-dependent parameter
spikes the position near periapsis while leaving the window-averaged error
almost unchanged; bounding the worst case keeps the realized quantized error
under control everywhere.

\begin{lstlisting}[language=cpp, caption={Worst-case perturbation used by the per-parameter bit search, \code{cpp/lupnt/applications/ephemeris/ephemeris\_app.cc}.}, label={lst:eph-bits}]
// cpp/lupnt/applications/ephemeris/ephemeris_app.cc :: SearchParamBits

const MatXd baseline = eph.Eval(t_s, params);
auto perturbed_error = [&](int k) {
  const double eps = std::pow(2.0, -k);
  VecXd p = params;
  p(idx) += eps;
  const MatXd plus = eph.Eval(t_s, p);
  p(idx) -= 2.0 * eps;
  const MatXd minus = eph.Eval(t_s, p);
  const double err_plus  = (plus.leftCols(3)  - baseline.leftCols(3)).rowwise().norm().maxCoeff();
  const double err_minus = (minus.leftCols(3) - baseline.leftCols(3)).rowwise().norm().maxCoeff();
  return std::max(err_plus, err_minus);
};
int lo = 0, hi = max_bits;                    // max_bits = 50
if (perturbed_error(lo) <= precision_m) return lo;
if (perturbed_error(hi) > precision_m) return hi;
while (hi - lo > 1) {
  const int mid = (lo + hi) / 2;
  if (perturbed_error(mid) <= precision_m) { hi = mid; } else { lo = mid; }
}
return hi;
\end{lstlisting}

The search is bracketed on $[0, 50]$ and returns the endpoint if the tolerance
is already met at $k = 0$ or still unmet at $k = 50$. The bisection is valid
because the perturbation magnitude, and hence the induced position offset, is
monotone in $k$.

\subsection{Range, sign, and margin bits}

Each parameter's field width is the sum of an integer part sized from its
dynamic range, the fractional part from \cref{eq:eph-bits-k}, a sign bit
$b_s \in \{0,1\}$, and a margin bit $b_m \in \{0,1\}$
\citep[Eqs.~7.11--7.12]{iiyama2026thesis}:
\begin{equation}
  \label{eq:eph-bits-range}
  b_j = \max\!\Big(\big\lceil \log_2\!\big(\alpha_{j,\max} - \alpha_{j,\min}\big)\big\rceil,\ 1\Big)
        + k_j + b_s + b_m,
  \qquad
  \text{Message size} = \sum_j b_j .
\end{equation}
The $\max(\cdot,1)$ guarantees at least one integer bit even for a parameter
whose range is below unity, and a non-positive range likewise falls back to one
bit.

\begin{lstlisting}[language=cpp, caption={Field-width rule, \code{cpp/lupnt/applications/ephemeris/ephemeris\_app.cc}.}, label={lst:eph-rangebits}]
// cpp/lupnt/applications/ephemeris/ephemeris_app.cc :: RangeBits

int RangeBits(double range, int bits_frac, bool signed_range, bool margin_bit) {
  const int bits_range
      = range > 0.0 ? std::max(1, static_cast<int>(std::ceil(std::log2(range)))) : 1;
  return bits_range + std::max(bits_frac, 0) + (signed_range ? 1 : 0) + (margin_bit ? 1 : 0);
}
\end{lstlisting}

The four parameter classes and their scale-factor conventions --- which follow
the GPS LNAV/CNAV scale-factor tables \citep{icdgps200} --- are given in
\cref{tab:eph-bitclasses}. Classification is by name prefix
(\cpp{IsAngleParam}, \cpp{IsEccentricityParam}, \cpp{IsSemiMajorAxisParam}),
tested in that order, so \code{argp} is classified as an angle even though it
also begins with \code{a}.

\begin{table}[H]
  \centering
  \caption{Bit-allocation classes. $\alpha_{j,\max}$ denotes the largest
    absolute value of parameter $j$ observed across all fitted windows.}
  \label{tab:eph-bitclasses}
  \begin{tabularx}{\textwidth}{@{}l l c c L@{}}
    \toprule
    Class & Range & $b_s$ & $b_m$ & Members (name prefix) \\
    \midrule
    Angle & $2\pi$ & $1$ & $0$ &
      \code{i}, \code{raan}, \code{argp}, \code{M}, \code{u} --- including all
      their polynomial and Fourier coefficients \\
    Eccentricity & $1$ & $0$ & $0$ & \code{e} and its coefficients \\
    Semi-major axis & $\alpha_{j,\max}$ & $0$ & $1$ &
      \code{a}, \code{a\_ref} and the $a$ coefficients \\
    Other & $\alpha_{j,\max}$ & $1$ & $1$ &
      Cartesian Chebyshev/Fourier coefficients \code{x\_*}, \code{y\_*},
      \code{z\_*} \\
    \bottomrule
  \end{tabularx}
\end{table}

Assigning the full $2\pi$ range to every coefficient of an angular element is
conservative for the higher-order coefficients, which are typically several
orders of magnitude smaller; it trades a few bits for a classification rule
that needs no per-coefficient tuning.

Indices $0$ and $1$ --- $t_\mathrm{ref}$ and $t_\mathrm{fit}$ --- are excluded
from the budget entirely. They are frame/timing metadata carried by the message
header rather than quantized model parameters.

\subsection{Sweep and realized error}

\cpp{SummarizeWindows} runs two passes over the \code{num\_windows} sampled
windows of a given length:
\begin{enumerate}[leftmargin=1.6em]
  \item \textbf{Sizing.} Fit every window; for each parameter take the maximum
    over windows of both $k_j$ from \cref{eq:eph-bits-k} and
    $|\alpha_j|$. The message must be sized for the worst window, not the
    average one.
  \item \textbf{Realization.} Re-quantize each window's parameters onto the
    chosen grid, $\hat\alpha_j = \operatorname{round}(\alpha_j/2^{-k_j})\,2^{-k_j}$,
    and evaluate \cpp{EvalError} with the \emph{quantized} vector. The
    reported RMS and 95th-percentile errors are therefore what a receiver
    decoding the actual broadcast bits would obtain, consistent with the
    reported \code{total\_bits}.
\end{enumerate}

The default tolerance is $\delta_x = \SI{0.01}{\meter}$
(\code{datasize\_precision\_m} in \cpp{EphemerisSimulationConfig}); the shipped
study configuration \file{configs/ephemeris.yaml} relaxes it to
\SI{0.05}{\meter} and sweeps
\code{fit\_window\_minutes: [30, 60, 120]} over four windows of an ELFO arc.

\begin{lupntnote}
  \textbf{Implementation vs.\ thesis.} Thesis Eqs.~7.9--7.10 size bits from a
  closed-form sensitivity,
  $k_j = -\big\lceil\log_2\!\big(\delta_x / |\partial f/\partial\alpha_j|\big)\big\rceil$,
  followed by a $\pm 1$ refinement. \lupnt{} reaches the same target with a
  direct binary search over the realized worst-case quantized error, which
  needs no derivative of the model and automatically accounts for the
  nonlinearity of the reconstruction near periapsis. The default tolerance is
  $\delta_x = \SI{0.01}{\meter}$. The LunaNet Interoperability Specification
  \citep{lunanet2022} allots roughly \textbf{\SI{900}{\bit}} to ephemeris data
  per broadcast frame \citep[\S1.2.3]{iiyama2026thesis}; the results
  \citep[Table~7.2]{iiyama2026thesis} meet the $3\sigma$ SISE thresholds
  (\SI{13.43}{\meter} and \SI{1.2}{\milli\meter\per\second} for LCRNS) within
  that budget for ELFO fits up to about \SI{240}{\minute}, with Chebyshev plus
  osculating elements --- plus Fourier terms for windows
  $\ge \SI{240}{\minute}$ --- as the best accuracy-per-bit representation. The
  same conclusion is reached independently in
  \citep{cortinovis2023ephemeris,cortinovis2024navi}.
\end{lupntnote}

% ===========================================================================
\section{Broadcast-Message Generation}
\label{sec:eph-broadcast}

\cpp{EphemerisGenApp}
(\file{cpp/lupnt/applications/ephemeris/ephemeris_gen_app.h},
\file{ephemeris_gen_app.cc}) is the \cpp{LunaNetSubApp} that produces both
message types from a satellite's own predicted arc, as a real vehicle would.
It is attached to a \cpp{LunaNetSatApp} with \cpp{AddSubApp}.

On each \cpp{Step(t)} it checks whether either message is due for refresh. When
one is, \cpp{GenerateFromAgent}:
\begin{enumerate}[leftmargin=1.6em]
  \item samples the owning agent's \emph{predicted} state at
    \code{n\_samples} epochs spaced uniformly over
    $[t_\mathrm{gen},\, t_\mathrm{gen} + T_\mathrm{window}]$ via
    \cpp{Agent::GetStateAt} --- i.e.\ the satellite propagates its own orbit
    forward before broadcasting;
  \item converts the sampled arc from the agent's dynamics frame to
    \code{output\_frame} with \cpp{ConvertFrame} if they differ (typically
    MOON\_CI $\to$ MOON\_PA);
  \item fits the corresponding model and appends a \cpp{BroadcastMessage}
    holding $[t_\mathrm{generated}, t_\mathrm{start}, t_\mathrm{end},
    \mathcal{F}, \vec{\alpha}]$.
\end{enumerate}
\cpp{LatestEphemeris(t)} and \cpp{LatestAlmanac(t)} return the last message
whose validity window contains $t$, mirroring a receiver selecting the
currently valid broadcast page; later messages supersede earlier ones.

\begin{table}[H]
  \centering
  \caption{\cpp{EphemerisGenConfig} defaults, \code{cpp/lupnt/applications/ephemeris/ephemeris\_gen\_app.h}. A non-positive refresh interval means ``refresh once per validity window''.}
  \label{tab:eph-genconfig}
  \begin{tabularx}{\textwidth}{@{}l l L@{}}
    \toprule
    Field & Default & Meaning \\
    \midrule
    \code{ephemeris\_window\_s} & \SI{2}{\hour} & Precise-ephemeris validity window \\
    \code{ephemeris\_refresh\_s} & $0$ & Refresh cadence; $\le 0$ means once per window \\
    \code{ephemeris\_fit\_samples} & $121$ & Uniform samples of the predicted arc \\
    \code{almanac\_window\_s} & \SI{15}{\dayunit} & Coarse-almanac validity window \\
    \code{almanac\_refresh\_s} & $0$ & Refresh cadence \\
    \code{almanac\_fit\_samples} & $361$ & Uniform samples of the predicted arc \\
    \code{output\_frame} & \code{Frame::MOON\_CI} &
      Broadcast frame; \cpp{Setup} forces both models' \code{frame} to it \\
    \bottomrule
  \end{tabularx}
\end{table}

\begin{lstlisting}[language=cpp, caption={Predicted-arc sampling and frame conversion before fitting, \code{cpp/lupnt/applications/ephemeris/ephemeris\_gen\_app.cc}.}, label={lst:eph-genapp}]
// cpp/lupnt/applications/ephemeris/ephemeris_gen_app.cc :: GenerateFromAgent

// Sample the agent's predicted state over the validity window, i.e. propagate
// the satellite's own orbit forward the way a real vehicle would before
// broadcasting an ephemeris/almanac for the upcoming window.
VecXd t_s(n_samples);
MatXd rv(n_samples, 6);
Frame agent_frame = Frame::MOON_CI;
for (int i = 0; i < n_samples; ++i) {
  const double tau = t_gen_s + window_s * static_cast<double>(i) / (n_samples - 1);
  Cart6 x = agent->GetStateAt(Real(tau));
  if (i == 0) agent_frame = x.GetFrame();
  t_s(i) = tau;
  rv.row(i) = x.cast<double>().transpose();
}

// Broadcast in output_frame: convert the sampled arc from the agent's dynamics
// frame if needed (tau is the absolute epoch each state was propagated to).
if (agent_frame != config_.output_frame) {
  for (int i = 0; i < n_samples; ++i) {
    const Vec6 rv_in = rv.row(i).transpose().cast<Real>();
    const Vec6 rv_out = ConvertFrame(Real(t_s(i)), rv_in, agent_frame, config_.output_frame);
    rv.row(i) = rv_out.cast<double>().transpose();
  }
}
\end{lstlisting}

Note the two different roles of the ``window'' in the two applications.
\cpp{EphemerisApp} (\cref{sec:eph-bits}) fits \emph{past} truth samples to
characterise the accuracy/datasize trade-off; \cpp{EphemerisGenApp} fits a
\emph{forward-predicted} arc, so its realized user error also contains the
satellite's own orbit-prediction error, which the fit statistics of
\cref{sec:eph-contract} do not capture.

% ===========================================================================
\section{Model Boundaries}
\label{sec:eph-boundaries}

\begin{itemize}[leftmargin=1.4em]
  \item \textbf{Frame is the caller's responsibility.} Inputs to \cpp{Fit} and
    \cpp{EvalError} must already be expressed in \code{options.frame};
    broadcasting in MOON\_PA requires the caller to \cpp{ConvertFrame} from MCI
    first. Neither model performs a frame conversion --- it only applies the
    scalar $\vec{\omega}\times\vec{r}$ correction of \cref{eq:eph-spin}, which
    is correct for a pure $z$-axis spin and ignores lunar libration and pole
    motion within the window.
  \item \textbf{Fourier terms require the baseline.} \cpp{LansEphemeris}
    Fourier terms require \code{use\_keplerian\_baseline}, because the argument
    of latitude comes from the baseline orbit. The almanac always carries the
    baseline in the sense that osculating elements \emph{are} the
    representation.
  \item \textbf{Node placement is a caller responsibility.} Chebyshev--Lobatto
    versus uniform sampling is decided outside the models, which fit the
    supplied epochs directly. Uniform sampling
    (\code{ephemeris\_fit\_samples = 121}) is what
    \cpp{EphemerisGenApp} uses, so Runge behaviour at the window edges is not
    suppressed by node placement and instead has to be paid for in polynomial
    order.
  \item \textbf{Velocity is analytic throughout.} Almanac velocity is the
    analytic two-body velocity, not a finite difference; ephemeris velocity is
    the analytic derivative of the Chebyshev and Fourier bases plus the
    baseline's rotated velocity. Consequently the broadcast velocity is
    exactly the derivative of the broadcast position in the ephemeris case,
    and is exactly two-body-consistent with the reconstructed elements in the
    almanac case.
  \item \textbf{The mean-anomaly integral is grid-dependent.} The almanac's
    $M_\mathrm{nom}$ is a cumulative trapezoid quadrature
    \cref{eq:eph-alm-cumtrapz} evaluated on the query grid. Two receivers
    evaluating the same message on different epoch grids obtain slightly
    different $M_\mathrm{nom}$; the difference is second order in the grid
    spacing and negligible at the almanac's accuracy level, but it means the
    reconstruction is not strictly a pointwise closed-form function of $t$.
  \item \textbf{The bit budget lives outside the models.} The
    quantization analysis is implemented in \cpp{EphemerisApp}, not in
    \cpp{LansEphemeris} or \cpp{LansAlmanac}, and \cpp{EphemerisGenApp} stores
    \emph{unquantized} double-precision parameters. A simulation that needs
    the realized broadcast error must apply the quantization itself.
  \item \textbf{Expected accuracy.} Over ELFO fit windows up to about
    \SI{240}{\minute}, the ephemeris model meets the LCRNS $3\sigma$ SISE
    thresholds of \SI{13.43}{\meter} and
    \SI{1.2}{\milli\meter\per\second} inside the LunaNet ephemeris allotment
    of roughly \SI{900}{\bit} \citep{lunanet2022,iiyama2025ephemeris}. The
    almanac is a deliberately coarse, long-validity model --- accurate enough
    for acquisition and visibility prediction over a 15-day window, and orders
    of magnitude less accurate than the ephemeris. It is not a fallback
    ranging source.
\end{itemize}
